\documentclass[11pt,a4paper]{article}
\usepackage[utf8]{inputenc}
\usepackage[T1]{fontenc}
\usepackage[catalan]{babel}
\usepackage{lmodern}
\usepackage[sfdefault,lining]{FiraSans}
\usepackage[a4paper,top=2.3cm,bottom=2.2cm,left=1.8cm,right=1.8cm,headheight=1.3cm,headsep=0.4cm,footskip=1cm]{geometry}
\usepackage{amsmath,amssymb}
\usepackage{icomma}
\usepackage{graphicx}
\usepackage[table]{xcolor}
\usepackage{tikz}
\usepackage{fontawesome5}
\usepackage{ccicons}
\usepackage{fancyhdr}
\usepackage{lastpage}
\usepackage{tasks}
\usepackage{needspace}
\usepackage[most]{tcolorbox}
\usepackage{array}
\usepackage[hidelinks,pdftitle={Fitxa d'exercicis de la unitat 3: Equacions i inequacions. 4t d'ESO Matemàtiques B},pdfsubject={Matemàtiques B 4t ESO},pdfkeywords={equacions de segon grau, biquadrades, Ruffini, factoritzades, fraccions algebraiques, radicals, exponencials, logarítmiques, inequacions}]{hyperref}

\definecolor{navy}{HTML}{1B2748}
\definecolor{accent}{HTML}{E8603A}
\definecolor{teal}{HTML}{1F9E90}
\definecolor{soft}{HTML}{EEF1F8}
\definecolor{muted}{HTML}{5D6685}

\setlength{\parindent}{0pt}
\setlength{\parskip}{3pt}
\settasks{label=\alph*), label-format=\bfseries\color{navy}, label-width=1.3em, item-indent=1.8em, column-sep=0.6em, after-item-skip=5pt, before-skip=3pt, after-skip=4pt}

% ---------- capçalera i peu ----------
\newcommand{\llicencia}{\href{https://creativecommons.org/licenses/by-nc-sa/4.0/deed.ca}{\ccbyncsa}\ \ Llicència Creative Commons Reconeixement-NoComercial-CompartirIgual 4.0 Internacional (CC BY-NC-SA 4.0)}
\pagestyle{fancy}
\fancyhf{}
\fancyhead[L]{\raisebox{-0.35\height}{\includegraphics[height=0.95cm]{logo-gva.png}}}
\fancyhead[C]{\small\color{muted}Fitxa d'exercicis · Unitat 3: Equacions i inequacions · 4t ESO B}
\fancyhead[R]{\raisebox{-0.35\height}{\includegraphics[height=0.95cm]{logo-vermellar.png}}}
\fancyfoot[L]{\scriptsize\color{muted}\llicencia}
\fancyfoot[R]{\small\color{muted}\thepage}
\renewcommand{\headrulewidth}{0.4pt}
\renewcommand{\headrule}{\color{muted!40}\hrule height \headrulewidth}
\fancypagestyle{primera}{\fancyhf{}\renewcommand{\headrulewidth}{0pt}
  \fancyfoot[L]{\scriptsize\color{muted}\llicencia}
  \fancyfoot[R]{\small\color{muted}\thepage}}

% ---------- nivells i exercicis ----------
\newcommand{\estrelles}[1]{{\color{accent}\ifcase#1\or\faStar\faStar[regular]\faStar[regular]\or\faStar\faStar\faStar[regular]\or\faStar\faStar\faStar\fi}}
\newcounter{exer}
\newcommand{\ex}[1]{\par\needspace{4\baselineskip}\medskip\refstepcounter{exer}\noindent{\color{navy}\bfseries\theexer.}\ \ {\small\estrelles{#1}}\ \ }

\newcommand{\seccio}[2]{\par\needspace{16\baselineskip}\bigskip
  \begin{tcolorbox}[colback=navy,colframe=navy,coltext=white,arc=2mm,boxrule=0pt,left=3mm,right=3mm,top=1.2mm,bottom=1.2mm]
  {\large\bfseries #1\quad #2}
  \end{tcolorbox}\vspace{-2pt}}

\newcommand{\exf}[1]{\par\medskip\refstepcounter{exer}\noindent{\color{navy}\bfseries\theexer.}\ \ {\small\estrelles{#1}}\ \ }
\newcolumntype{C}[1]{>{\centering\arraybackslash}p{#1}}
% quadre de recordatori (només a la fitxa d'exercicis)
\newcommand{\recorda}[1]{\par\smallskip\begin{tcolorbox}[colback=accent!7,colframe=accent!55,arc=2mm,boxrule=0.5pt,left=2.5mm,right=2.5mm,top=1mm,bottom=1mm]\small{\color{accent}\faLightbulb[regular]}\ \textbf{Recorda:}\ #1\end{tcolorbox}}
\newcommand{\res}[1]{#1}\newcommand{\rest}[1]{#1}

\def\fmneed{11}% que el títol d'una secció no quede sol al final del full
% =====================================================================
% Capçalera FentMates de les fitxes: un color per a cada curs.
% Els colors són els de la pàgina inicial (LEVELS a assets/js/materials.js);
% el logo ve de _eines/logo/pdf/ (python3 _eines/logo/genera.py).
%
% Ús, al final del preàmbul de la unitat:
%   % =====================================================================
% Capçalera FentMates de les fitxes: un color per a cada curs.
% Els colors són els de la pàgina inicial (LEVELS a assets/js/materials.js);
% el logo ve de _eines/logo/pdf/ (python3 _eines/logo/genera.py).
%
% Ús, al final del preàmbul de la unitat:
%   % =====================================================================
% Capçalera FentMates de les fitxes: un color per a cada curs.
% Els colors són els de la pàgina inicial (LEVELS a assets/js/materials.js);
% el logo ve de _eines/logo/pdf/ (python3 _eines/logo/genera.py).
%
% Ús, al final del preàmbul de la unitat:
%   \input{../comu/fentmates.tex}
%   \fmcurs{1eso}
% i, al començament del document:
%   \capcalera{Matemàtiques}{1r ESO}{Fitxa d'exercicis · Unitat 2}{Divisibilitat}
% Opcional abans de \input: \def\fmneed{14} (espai mínim abans de cada secció).
% =====================================================================
\makeatletter
\def\fm@logodir{../../logo/pdf}
\providecommand{\fmneed}{10}
\def\fm@cursnom{}

% clau del curs, color i color del text damunt del color
\newcommand{\fm@curs}[3]{\@namedef{fm@color@#1}{#2}\@namedef{fm@text@#1}{#3}}
\fm@curs{1eso}{E6B84A}{navy}
\fm@curs{2eso}{3B82F6}{white}
\fm@curs{3eso}{FF7A50}{white}
\fm@curs{4eso-a}{A78BFA}{white}
\fm@curs{4eso-b}{6D28D9}{white}
\fm@curs{1batx-mi}{9B1C1C}{white}
\fm@curs{1batx-ccss}{EF6B6B}{white}
\fm@curs{2batx-mii}{1E6B47}{white}
\fm@curs{2batx-ccss}{2BB3A3}{white}

\newcommand{\fmcurs}[1]{%
  \@ifundefined{fm@color@#1}{\PackageError{fentmates}{Curs desconegut: #1}{}}{}%
  \edef\fm@tmp{\noexpand\definecolor{curs}{HTML}{\@nameuse{fm@color@#1}}%
               \noexpand\colorlet{curstext}{\@nameuse{fm@text@#1}}}\fm@tmp
  \colorlet{cursfosc}{curs!72!black}%
  \def\fm@logo{\fm@logodir/vertical-#1.pdf}%
  \def\fm@marca{\fm@logodir/marca-#1.pdf}}

% ---------- capçalera de la primera pàgina ----------
% #1 matèria (petit), #2 curs, #3 tipus de document i unitat, #4 títol
\newcommand{\capcalera}[4]{%
  \gdef\fm@cursnom{#2}%
  \par\noindent
  \begin{tcolorbox}[colback=curs,colframe=curs,boxrule=0pt,arc=2.2mm,width=3.5cm,height=2.15cm,
      valign=center,left=3mm,right=2mm,top=1mm,bottom=1mm,nobeforeafter,box align=center]
    \color{curstext}{\footnotesize #1}\par\vspace{1pt}{\fontsize{19}{21}\selectfont\bfseries #2\par}%
  \end{tcolorbox}\hspace{5mm}%
  \begin{minipage}[c]{\dimexpr\linewidth-3.5cm-5mm-2.9cm\relax}
    {\color{cursfosc}\small\bfseries\MakeUppercase{#3}}\par\vspace{2pt}
    {\color{navy}\fontsize{20}{23}\selectfont\bfseries #4\par}
  \end{minipage}\hfill
  \begin{minipage}[c]{2.6cm}\raggedleft\includegraphics[width=2.6cm]{\fm@logo}\end{minipage}\par
  \vspace{2pt}}

% ---------- barres de secció del color del curs ----------
\renewcommand{\seccio}[2]{\par\needspace{\fmneed\baselineskip}\bigskip
  \begin{tcolorbox}[colback=curs,colframe=curs,coltext=curstext,arc=2mm,boxrule=0pt,left=3mm,right=3mm,top=1.2mm,bottom=1.2mm]
  {\large\bfseries #1\quad #2}
  \end{tcolorbox}\vspace{-2pt}}

% ---------- estrelles de dificultat del color del curs (a Batxillerat no se'n posen) ----------
\@ifundefined{estrelles}{}{%
  \renewcommand{\estrelles}[1]{{\color{curs}\ifcase#1\or\faStar\faStar[regular]\faStar[regular]\or\faStar\faStar\faStar[regular]\or\faStar\faStar\faStar\fi}}}

% ---------- capçalera i peu de les pàgines ----------
% primera pàgina: logos institucionals al peu
\fancypagestyle{primera}{\fancyhf{}\renewcommand{\headrulewidth}{0pt}
  \fancyfoot[L]{\raisebox{-0.32\height}{\includegraphics[height=0.78cm]{logo-gva.png}}\hspace{5mm}%
    \raisebox{-0.32\height}{\includegraphics[height=0.78cm]{logo-vermellar.png}}}
  \fancyfoot[R]{\scriptsize\color{muted}\href{https://creativecommons.org/licenses/by-nc-sa/4.0/deed.ca}{\ccbyncsa}\ \ CC BY-NC-SA 4.0 · fentmates.org\qquad\small\thepage}}
% la resta: el símbol a l'esquerra i el curs a la dreta, amb el títol al centre com abans
\fancyhead[L]{\raisebox{-0.3\height}{\includegraphics[height=0.62cm]{\fm@marca}}}
\fancyhead[R]{\small\bfseries\color{cursfosc}\fm@cursnom}
\renewcommand{\headrule}{\color{curs!45}\hrule height \headrulewidth}
\makeatother

%   \fmcurs{1eso}
% i, al començament del document:
%   \capcalera{Matemàtiques}{1r ESO}{Fitxa d'exercicis · Unitat 2}{Divisibilitat}
% Opcional abans de \input: \def\fmneed{14} (espai mínim abans de cada secció).
% =====================================================================
\makeatletter
\def\fm@logodir{../../logo/pdf}
\providecommand{\fmneed}{10}
\def\fm@cursnom{}

% clau del curs, color i color del text damunt del color
\newcommand{\fm@curs}[3]{\@namedef{fm@color@#1}{#2}\@namedef{fm@text@#1}{#3}}
\fm@curs{1eso}{E6B84A}{navy}
\fm@curs{2eso}{3B82F6}{white}
\fm@curs{3eso}{FF7A50}{white}
\fm@curs{4eso-a}{A78BFA}{white}
\fm@curs{4eso-b}{6D28D9}{white}
\fm@curs{1batx-mi}{9B1C1C}{white}
\fm@curs{1batx-ccss}{EF6B6B}{white}
\fm@curs{2batx-mii}{1E6B47}{white}
\fm@curs{2batx-ccss}{2BB3A3}{white}

\newcommand{\fmcurs}[1]{%
  \@ifundefined{fm@color@#1}{\PackageError{fentmates}{Curs desconegut: #1}{}}{}%
  \edef\fm@tmp{\noexpand\definecolor{curs}{HTML}{\@nameuse{fm@color@#1}}%
               \noexpand\colorlet{curstext}{\@nameuse{fm@text@#1}}}\fm@tmp
  \colorlet{cursfosc}{curs!72!black}%
  \def\fm@logo{\fm@logodir/vertical-#1.pdf}%
  \def\fm@marca{\fm@logodir/marca-#1.pdf}}

% ---------- capçalera de la primera pàgina ----------
% #1 matèria (petit), #2 curs, #3 tipus de document i unitat, #4 títol
\newcommand{\capcalera}[4]{%
  \gdef\fm@cursnom{#2}%
  \par\noindent
  \begin{tcolorbox}[colback=curs,colframe=curs,boxrule=0pt,arc=2.2mm,width=3.5cm,height=2.15cm,
      valign=center,left=3mm,right=2mm,top=1mm,bottom=1mm,nobeforeafter,box align=center]
    \color{curstext}{\footnotesize #1}\par\vspace{1pt}{\fontsize{19}{21}\selectfont\bfseries #2\par}%
  \end{tcolorbox}\hspace{5mm}%
  \begin{minipage}[c]{\dimexpr\linewidth-3.5cm-5mm-2.9cm\relax}
    {\color{cursfosc}\small\bfseries\MakeUppercase{#3}}\par\vspace{2pt}
    {\color{navy}\fontsize{20}{23}\selectfont\bfseries #4\par}
  \end{minipage}\hfill
  \begin{minipage}[c]{2.6cm}\raggedleft\includegraphics[width=2.6cm]{\fm@logo}\end{minipage}\par
  \vspace{2pt}}

% ---------- barres de secció del color del curs ----------
\renewcommand{\seccio}[2]{\par\needspace{\fmneed\baselineskip}\bigskip
  \begin{tcolorbox}[colback=curs,colframe=curs,coltext=curstext,arc=2mm,boxrule=0pt,left=3mm,right=3mm,top=1.2mm,bottom=1.2mm]
  {\large\bfseries #1\quad #2}
  \end{tcolorbox}\vspace{-2pt}}

% ---------- estrelles de dificultat del color del curs (a Batxillerat no se'n posen) ----------
\@ifundefined{estrelles}{}{%
  \renewcommand{\estrelles}[1]{{\color{curs}\ifcase#1\or\faStar\faStar[regular]\faStar[regular]\or\faStar\faStar\faStar[regular]\or\faStar\faStar\faStar\fi}}}

% ---------- capçalera i peu de les pàgines ----------
% primera pàgina: logos institucionals al peu
\fancypagestyle{primera}{\fancyhf{}\renewcommand{\headrulewidth}{0pt}
  \fancyfoot[L]{\raisebox{-0.32\height}{\includegraphics[height=0.78cm]{logo-gva.png}}\hspace{5mm}%
    \raisebox{-0.32\height}{\includegraphics[height=0.78cm]{logo-vermellar.png}}}
  \fancyfoot[R]{\scriptsize\color{muted}\href{https://creativecommons.org/licenses/by-nc-sa/4.0/deed.ca}{\ccbyncsa}\ \ CC BY-NC-SA 4.0 · fentmates.org\qquad\small\thepage}}
% la resta: el símbol a l'esquerra i el curs a la dreta, amb el títol al centre com abans
\fancyhead[L]{\raisebox{-0.3\height}{\includegraphics[height=0.62cm]{\fm@marca}}}
\fancyhead[R]{\small\bfseries\color{cursfosc}\fm@cursnom}
\renewcommand{\headrule}{\color{curs!45}\hrule height \headrulewidth}
\makeatother

%   \fmcurs{1eso}
% i, al començament del document:
%   \capcalera{Matemàtiques}{1r ESO}{Fitxa d'exercicis · Unitat 2}{Divisibilitat}
% Opcional abans de \input: \def\fmneed{14} (espai mínim abans de cada secció).
% =====================================================================
\makeatletter
\def\fm@logodir{../../logo/pdf}
\providecommand{\fmneed}{10}
\def\fm@cursnom{}

% clau del curs, color i color del text damunt del color
\newcommand{\fm@curs}[3]{\@namedef{fm@color@#1}{#2}\@namedef{fm@text@#1}{#3}}
\fm@curs{1eso}{E6B84A}{navy}
\fm@curs{2eso}{3B82F6}{white}
\fm@curs{3eso}{FF7A50}{white}
\fm@curs{4eso-a}{A78BFA}{white}
\fm@curs{4eso-b}{6D28D9}{white}
\fm@curs{1batx-mi}{9B1C1C}{white}
\fm@curs{1batx-ccss}{EF6B6B}{white}
\fm@curs{2batx-mii}{1E6B47}{white}
\fm@curs{2batx-ccss}{2BB3A3}{white}

\newcommand{\fmcurs}[1]{%
  \@ifundefined{fm@color@#1}{\PackageError{fentmates}{Curs desconegut: #1}{}}{}%
  \edef\fm@tmp{\noexpand\definecolor{curs}{HTML}{\@nameuse{fm@color@#1}}%
               \noexpand\colorlet{curstext}{\@nameuse{fm@text@#1}}}\fm@tmp
  \colorlet{cursfosc}{curs!72!black}%
  \def\fm@logo{\fm@logodir/vertical-#1.pdf}%
  \def\fm@marca{\fm@logodir/marca-#1.pdf}}

% ---------- capçalera de la primera pàgina ----------
% #1 matèria (petit), #2 curs, #3 tipus de document i unitat, #4 títol
\newcommand{\capcalera}[4]{%
  \gdef\fm@cursnom{#2}%
  \par\noindent
  \begin{tcolorbox}[colback=curs,colframe=curs,boxrule=0pt,arc=2.2mm,width=3.5cm,height=2.15cm,
      valign=center,left=3mm,right=2mm,top=1mm,bottom=1mm,nobeforeafter,box align=center]
    \color{curstext}{\footnotesize #1}\par\vspace{1pt}{\fontsize{19}{21}\selectfont\bfseries #2\par}%
  \end{tcolorbox}\hspace{5mm}%
  \begin{minipage}[c]{\dimexpr\linewidth-3.5cm-5mm-2.9cm\relax}
    {\color{cursfosc}\small\bfseries\MakeUppercase{#3}}\par\vspace{2pt}
    {\color{navy}\fontsize{20}{23}\selectfont\bfseries #4\par}
  \end{minipage}\hfill
  \begin{minipage}[c]{2.6cm}\raggedleft\includegraphics[width=2.6cm]{\fm@logo}\end{minipage}\par
  \vspace{2pt}}

% ---------- barres de secció del color del curs ----------
\renewcommand{\seccio}[2]{\par\needspace{\fmneed\baselineskip}\bigskip
  \begin{tcolorbox}[colback=curs,colframe=curs,coltext=curstext,arc=2mm,boxrule=0pt,left=3mm,right=3mm,top=1.2mm,bottom=1.2mm]
  {\large\bfseries #1\quad #2}
  \end{tcolorbox}\vspace{-2pt}}

% ---------- estrelles de dificultat del color del curs (a Batxillerat no se'n posen) ----------
\@ifundefined{estrelles}{}{%
  \renewcommand{\estrelles}[1]{{\color{curs}\ifcase#1\or\faStar\faStar[regular]\faStar[regular]\or\faStar\faStar\faStar[regular]\or\faStar\faStar\faStar\fi}}}

% ---------- capçalera i peu de les pàgines ----------
% primera pàgina: logos institucionals al peu
\fancypagestyle{primera}{\fancyhf{}\renewcommand{\headrulewidth}{0pt}
  \fancyfoot[L]{\raisebox{-0.32\height}{\includegraphics[height=0.78cm]{logo-gva.png}}\hspace{5mm}%
    \raisebox{-0.32\height}{\includegraphics[height=0.78cm]{logo-vermellar.png}}}
  \fancyfoot[R]{\scriptsize\color{muted}\href{https://creativecommons.org/licenses/by-nc-sa/4.0/deed.ca}{\ccbyncsa}\ \ CC BY-NC-SA 4.0 · fentmates.org\qquad\small\thepage}}
% la resta: el símbol a l'esquerra i el curs a la dreta, amb el títol al centre com abans
\fancyhead[L]{\raisebox{-0.3\height}{\includegraphics[height=0.62cm]{\fm@marca}}}
\fancyhead[R]{\small\bfseries\color{cursfosc}\fm@cursnom}
\renewcommand{\headrule}{\color{curs!45}\hrule height \headrulewidth}
\makeatother

\fmcurs{4eso-b}

\begin{document}
\thispagestyle{primera}

\capcalera{Matemàtiques B}{4t ESO}{Fitxa d'exercicis · Unitat 3}{\raggedright Equacions i inequacions}
\vspace{10pt}
Nom i cognoms: \rule{8.4cm}{0.4pt}\hfill Grup: \rule{1.6cm}{0.4pt}\hfill Data: \rule{2.2cm}{0.4pt}
\vspace{8pt}

\begin{tcolorbox}[colback=soft,colframe=soft,arc=2mm,left=3mm,right=3mm,top=1.5mm,bottom=1.5mm]
\small Exercicis graduats per nivells de dificultat:\quad \estrelles{1}\ bàsic \qquad \estrelles{2}\ mitjà \qquad \estrelles{3}\ avançat.\\[2pt] Resol-los al quadern mostrant tots els passos. Comprova sempre les solucions de les equacions amb fraccions, amb radicals i logarítmiques, i dona les solucions de les inequacions com a interval i representades en la recta.
\end{tcolorbox}

% =====================================================
\seccio{1}{Equacions de segon grau}

\exf{1} Resol les equacions incompletes:\begin{tasks}(3)
  \task $x^2-16=0$
  \task $3x^2+12x=0$
  \task $2x^2-50=0$
  \task $x^2+9=0$
  \task $5x^2=3x$
  \task $4x^2-1=0$
\end{tasks}

\ex{1} Resol les equacions:\begin{tasks}(3)
  \task $x^2-5x+6=0$
  \task $x^2+2x-15=0$
  \task $2x^2-7x+3=0$
  \task $x^2-6x+9=0$
  \task $x^2+x+3=0$
  \task $6x^2-x-1=0$
\end{tasks}\recorda{$ax^2+bx+c=0\Rightarrow x=\dfrac{-b\pm\sqrt{b^2-4ac}}{2a}$.}

\ex{2} Opera fins a obtindre la forma $ax^2+bx+c=0$ i resol:\begin{tasks}(2)
  \task $(x+2)^2=3x+10$
  \task $x(x-3)=2(x+7)$
  \task $(2x-1)(x+3)-x^2=3$
  \task $\dfrac{x^2}{3}-\dfrac x2=\dfrac32$
\end{tasks}

\ex{2} Sense resoldre-les, indica quantes solucions té cada equació:\begin{tasks}(3)
  \task $3x^2-2x+5=0$
  \task $4x^2-12x+9=0$
  \task $x^2+x-1=0$
\end{tasks}Troba després $k$ perquè $x^2-6x+k=0$ tinga una solució doble.

\ex{3} Escriu una equació de segon grau que tinga per solucions $x=3$ i $x=-5$.


% =====================================================
\seccio{2}{Equacions biquadrades}

\exf{1} Resol les equacions biquadrades:\begin{tasks}(2)
  \task $x^4-10x^2+9=0$
  \task $x^4-5x^2-36=0$
  \task $x^4+3x^2+2=0$
  \task $4x^4-17x^2+4=0$
\end{tasks}\recorda{Fes el canvi $t=x^2$. Cada solució $t>0$ dona $x=\pm\sqrt t$; les $t<0$ no donen solucions.}

\ex{2} Resol:\begin{tasks}(3)
  \task $x^4-8x^2=0$
  \task $(x^2-1)(x^2+2)=10$
  \task $x^4-16=0$
\end{tasks}


% =====================================================
\seccio{3}{Equacions de grau major que 2}

\exf{2} Resol amb la regla de Ruffini:\begin{tasks}(2)
  \task $x^3-6x^2+11x-6=0$
  \task $x^3+2x^2-x-2=0$
  \task $x^3-3x+2=0$
  \task $2x^3+x^2-8x-4=0$
\end{tasks}\recorda{Prova amb els divisors del terme independent fins a arribar a una equació de segon grau.}

\ex{2} Passa tots els termes a un membre, factoritza i resol:\begin{tasks}(3)
  \task $x^4-x^3-7x^2+x+6=0$
  \task $x^3=4x$
  \task $x^4+x^3=x^2+x$
\end{tasks}


% =====================================================
\seccio{4}{Equacions factoritzades}

\exf{1} Resol sense multiplicar els factors:\begin{tasks}(2)
  \task $(x-4)(x+1)=0$
  \task $x(2x+3)(x-5)=0$
  \task $(x-1)^2(x+7)=0$
  \task $3(x+2)(4x-1)=0$
\end{tasks}\recorda{Un producte és zero si algun dels factors és zero.}

\ex{2} Resol:\begin{tasks}(2)
  \task $(x^2-16)(x+3)=0$
  \task $(x^2-3x+2)(x^2+5)=0$
  \task $x^3(x^2-2x-8)=0$
  \task $(x-2)(x+1)=4$
\end{tasks}


% =====================================================
\seccio{5}{Equacions amb fraccions algebraiques}

\exf{1} Resol:\begin{tasks}(2)
  \task $\dfrac5x-1=4$
  \task $\dfrac{x}{x-1}=\dfrac32$
  \task $\dfrac1x+\dfrac{2}{x+1}=1$
  \task $x+\dfrac4x=5$
\end{tasks}

\ex{2} Resol i comprova les solucions:\begin{tasks}(1)
  \task $\dfrac{x}{x+2}-\dfrac{1}{x-2}=\dfrac{8}{x^2-4}$
  \task $\dfrac{3}{x-1}-\dfrac{2}{x+1}=\dfrac{6}{x^2-1}$
\end{tasks}\recorda{Descarta les solucions que anul·len algun denominador.}

\ex{3} Resol $\dfrac{x+1}{x}-\dfrac{x}{x+1}=\dfrac56$.


% =====================================================
\seccio{6}{Equacions amb radicals}

\exf{1} Resol:\begin{tasks}(4)
  \task $\sqrt{x+5}=3$
  \task $\sqrt{2x-1}=5$
  \task $\sqrt{x^2+9}=5$
  \task $\sqrt{x+3}=-2$
\end{tasks}

\ex{2} Resol i comprova les solucions:\begin{tasks}(2)
  \task $\sqrt{x+2}=x$
  \task $x+\sqrt{x-1}=7$
  \task $\sqrt{3x+4}-x=0$
  \task $2\sqrt x=x-3$
\end{tasks}\recorda{Aïlla el radical, eleva al quadrat i comprova totes les solucions en l'equació inicial.}

\ex{3} Resol (hi ha dos radicals):\begin{tasks}(2)
  \task $\sqrt{x+7}-\sqrt x=1$
  \task $\sqrt{2x+1}+\sqrt{x-3}=4$
\end{tasks}


% =====================================================
\seccio{7}{Equacions exponencials}

\exf{1} Escriu els dos membres com a potències de la mateixa base i resol:\begin{tasks}(3)
  \task $2^x=32$
  \task $3^{x-2}=27$
  \task $5^{2x}=\frac{1}{25}$
  \task $4^x=8$
  \task $2^{x^2-1}=8$
  \task $10^{x+1}=0,001$
\end{tasks}

\ex{1} Resol prenent logaritmes (dona el resultat amb tres decimals):\begin{tasks}(4)
  \task $2^x=10$
  \task $3^x=50$
  \task $7^{x-1}=30$
  \task $1,05^x=2$
\end{tasks}\recorda{$a^x=b\Rightarrow x=\dfrac{\log b}{\log a}$.}

\ex{2} Resol traient factor comú:\begin{tasks}(3)
  \task $2^{x+1}+2^x=24$
  \task $3^{x+2}-3^x=72$
  \task $5^x+5^{x-1}=30$
\end{tasks}

\ex{3} Resol amb un canvi de variable:\begin{tasks}(3)
  \task $4^x-5\cdot2^x+4=0$
  \task $9^x-10\cdot3^x+9=0$
  \task $4^x-2^x-12=0$
\end{tasks}


% =====================================================
\seccio{8}{Equacions logarítmiques}

\exf{1} Resol aplicant la definició de logaritme:\begin{tasks}(4)
  \task $\log x=3$
  \task $\log_2(x+1)=4$
  \task $\ln x=0$
  \task $\log_5(3x-1)=2$
\end{tasks}

\ex{2} Resol i comprova les solucions:\begin{tasks}(2)
  \task $\log x+\log4=\log20$
  \task $\log(x+3)-\log(x-1)=\log3$
  \task $2\log x=\log(5x-6)$
  \task $\log x+\log(x+9)=1$
\end{tasks}\recorda{Agrupa en un sol logaritme a cada membre. Els arguments dels logaritmes han de ser positius.}

\ex{3} Resol:\begin{tasks}(2)
  \task $\log(x+6)=1+\log(x-3)$
  \task $\log_2x+\log_2(x-2)=3$
\end{tasks}


% =====================================================
\seccio{9}{Inequacions de primer grau}

\exf{1} Resol i representa la solució:\begin{tasks}(4)
  \task $2x-5>3$
  \task $3-x\le5$
  \task $4x+1<2x-7$
  \task $-3x\ge12$
\end{tasks}\recorda{Si multipliques o divideixes per un nombre negatiu, la desigualtat canvia de sentit.}

\ex{2} Resol:\begin{tasks}(2)
  \task $3(x-2)-2(x+1)\ge x-8$
  \task $\dfrac{x+1}{2}-\dfrac{x-3}{4}<2$
  \task $2x-\dfrac{3-x}{2}>\dfrac x3$
  \task $5x-1>5(x+1)$
\end{tasks}

\ex{2} Amb la tarifa A, un mòbil costa 40 € al mes. Amb la tarifa B costa 15 € al mes més 0,50 € per cada GB. A partir de quants GB és més cara la tarifa B?


% =====================================================
\seccio{10}{Inequacions de segon grau o superior}

\exf{2} Resol i representa la solució:\begin{tasks}(3)
  \task $x^2-9<0$
  \task $x^2-3x-10\ge0$
  \task $x^2+4x>0$
  \task $-x^2+2x+3>0$
  \task $x^2-4x+4\le0$
  \task $x^2+x+1>0$
\end{tasks}\recorda{Calcula les arrels, divideix la recta en intervals i estudia el signe en cada un.}

\ex{3} Resol:\begin{tasks}(3)
  \task $x^3-x>0$
  \task $(x-1)(x+2)(x-4)\le0$
  \task $x^3-2x^2-3x\ge0$
\end{tasks}

\ex{2} Els costats d'un rectangle fan $x$ i $x+3$ cm. Per a quins valors de $x$ l'àrea és menor que 40 cm$^2$?


% =====================================================
\seccio{11}{Problemes}

\exf{2} El producte de dos nombres naturals consecutius és 156. Quins són?

\ex{2} Un rectangle té 34 cm de perímetre i 13 cm de diagonal. Quant fan els costats?

\ex{2} Quants anys han de passar perquè 5000 € a un 3\,\% d'interés compost anual es convertisquen en 6000 €?

\ex{3} Una població de bacteris segueix la fórmula $N=500\cdot2^{t/3}$, amb $t$ en hores. Quan hi haurà 16\,000 bacteris?

\vfill
\begin{center}\color{muted}\small\faCheckCircle[regular]\ \ Ànim! Comprova els teus resultats amb les eines interactives dels apunts de la unitat.\end{center}

\end{document}
